% Generated by tools/edn_latex.py from reports/edn.md.
% Do not edit: change reports/edn.md and run `make edn`.
\ednentry{
  id = {EDN-66},
  title = {NFR-06 and NFR-09 are each split in two, and each ID keeps the half its citations mean},
  date = {2026-10-05},
  milestone = {M3: Quality Assurance},
  topic = {Requirements Engineering, Testing Strategy},
  participants = {@lukas2510, decided by the repository owner},
  decision = {NFR-06 and NFR-09 each promised two things that come due in different milestones and are verified by different evidence, so each is split in two, and every half has its own ID, acceptance criterion, verification route and milestone.
The existing ID is narrowed to the half its citations mean, and the other half gets the next free ID.
\textbf{NFR-06} keeps reproducibility: from a clean clone, \texttt{dvc pull} and \texttt{dvc repro} reproduce the splits and metrics within \ensuremath{\pm}0.1 percentage points; \textbf{[manual]}, the clean-clone drill of \href{https://github.com/mlops-2627q1-mds-upc/MLOPS_Recommenditos/blob/main/docs/docs/pipeline.md}{The DVC pipeline}, due M3.
\textbf{NFR-14} takes run provenance: every training run records its commit, its data version and its parameters; \textbf{[automated]}, by the tracking tests and a \texttt{train} test that reads a real run back, due M3.
\textbf{NFR-09} keeps the running system's security: a read-only public surface, the body-size limit, no TLS, the API key from the environment, non-root containers and the vulnerability scan; \textbf{[automated]}, due M4.
\textbf{NFR-15} takes credential hygiene: no secret in any commit of the history, a scan before every merge, and a leaked secret revoked rather than only removed; \textbf{[manual]}, the CI secret scan of EDN-71, due M3.},
  rationale = {\emph{Alternatives considered:}
\begin{itemize}[leftmargin=*, topsep=2pt]
\item \textbf{Option A (chosen): narrow the existing ID to one half, and give the other half a new ID.}
\newline Pros: on \texttt{main} before the split, 66 lines cited NFR-06, 15 of them \texttt{req} markers, and 18 cited NFR-09.
Apart from the requirement rows themselves and the passages that tell NFR-06's history, eleven of NFR-06's citations mean provenance -- the tracking seam and its tests, two documentation pages and EDN-54 -- and of NFR-09's only EDN-46's evidence pointer means secret hygiene.
Every other citation, the report's included, means the half the ID keeps and stays correct without an edit, and the twelve tests marked NFR-06 for determinism and the download pin keep their marker.
No ID is reused for a different requirement: a narrowed ID promises a subset of what it promised before and nothing new.
\newline Cons: a reader holding an older text can still read NFR-06 as including provenance, so the citations of the moved halves had to be found and repointed.
Earlier entries written under the combined meaning, EDN-45, EDN-55 and EDN-56, are records of their date and are not rewritten; EDN-56 in fact already argued the two halves of NFR-06 separately, which is the strongest sign they were two requirements.
\item \textbf{Option B: retire NFR-06 and NFR-09 and issue four new IDs.}
\newline Pros: no ID ever changes meaning, not even by narrowing.
\newline Cons: every one of those 84 citations becomes a pointer to a requirement that no longer exists, including three sections of the LaTeX report; the fifteen NFR-06 markers fail the per-test check until each is rewritten; and the matrix needs every requirement to have a table row, so a retired ID could only survive in prose, where nothing checks it.
\item \textbf{Option C: suffixes, NFR-06a and NFR-06b.}
\newline Pros: the family stays visible in the name.
\newline Cons: \texttt{tools/\allowbreak{}requirement\_\allowbreak{}matrix.\allowbreak{}py} accepts \texttt{(?:FR|NFR)-\allowbreak{}\textbackslash{}d\{2,\}} and nothing else, so the tool would have to change first; and a bare ``NFR-06'' in an existing citation would become ambiguous between the parent and its halves, which is the ambiguity the split exists to remove.
\item \textbf{Option D: keep one ID each, and let the ``Verified by'' cell name both routes.}
\newline Pros: no ID change at all.
\newline Cons: it is the state issue \#60 describes.
The parser requires exactly one of \textbf{[automated]} and \textbf{[manual]}, and the gate books an ID to one milestone, so NFR-09's hygiene half sat in M4 behind an API that does not exist yet while the token it protects was already on every contributor's disk, and NFR-06's tests could only ever make it ``named by a test''.
\end{itemize}
\par
\emph{Why:} The constraint is that a requirement ID is never reused for a different requirement, because the report and the EDN cite them; between A and B the question is only which citations break.
In A it is a bounded, listed set, the ones that mean the half that moved; in B it is all of them.
The halves get their milestones from when their evidence can exist, not from the requirement they came from: NFR-14 is due in M3 because the tracking seam (PR \#50) and the \texttt{train} stage's parameter logging (PR \#62) are both merged, and NFR-15 is due in M3 because the risk has existed since issue \#42 and PR \#50 put a DagsHub token into every contributor's \texttt{.\allowbreak{}env} and \texttt{.\allowbreak{}dvc/\allowbreak{}config.\allowbreak{}local} (EDN-46).
NFR-09 stays booked to M4 as decided with the split, although its image scan, its Dockerfile review and the check from outside the VM are M5 work; \texttt{tools/\allowbreak{}expected\_\allowbreak{}coverage.\allowbreak{}yaml} says so, so the delivery review notices if they are still missing then.
NFR-14 is \textbf{[automated]} because what a run records can be read back from a run; NFR-15 is \textbf{[manual]} because its evidence is a CI job and not a Pytest test, the route NFR-07 already uses for the CI run it names, and making it \textbf{[automated]} would have needed a test that only stands in for the scan.
Giving NFR-14 a route of its own also showed that it was not met: \texttt{train.\allowbreak{}num\_\allowbreak{}threads} is declared as a parameter of the \texttt{train} stage in \texttt{dvc.\allowbreak{}yaml} and changes the artefact (EDN-53), yet the run carried it only inside the \texttt{model.\allowbreak{}json} artefact, not as an MLflow parameter, so a comparison of runs by parameter could not see it, and the only test spot-checked two parameters.
So a test now reads the keys \texttt{dvc.\allowbreak{}yaml} declares for each \texttt{train@\allowbreak{}<variant>} out of \texttt{dvc.\allowbreak{}yaml} itself and compares them with the parameters of a real run; it failed on the code before the change, on \texttt{num\_\allowbreak{}threads} for all four variants, and passes since the stage logs it.
\texttt{train.\allowbreak{}mlflow\_\allowbreak{}experiment} is the one declared key recorded as what it is, the experiment the run belongs to, and the test checks it there rather than as a repeated parameter.
\texttt{train.\allowbreak{}py} and \texttt{recommenditos/\allowbreak{}tracking.\allowbreak{}py}, whose three comments also move from NFR-06 to NFR-14, are dependencies of the four \texttt{train} stages and of \texttt{evaluate}, so the change reruns them on the real snapshot.
The re-run on 2026-10-05 moved no number: \texttt{metrics.\allowbreak{}json}, the segment table and the masking sweep are unchanged, and inside each model bundle only \texttt{model.\allowbreak{}json} changed, by the provenance EDN-56 puts there, while every payload file kept its hash; the four new runs carry all of their stage's declared parameters, \texttt{num\_\allowbreak{}threads} among them.},
  ai = {Information seeking, Alternative generation, Alternative assessment, Recommendation, Solution generation},
  response = {Accepted with modifications},
  assessment = {The owner decided to split both requirements, which half comes due when, and that the hygiene half is verified by a secret scan; the issue and its comment are his.
How to split the IDs he delegated, under the constraint that no ID is reused.
AI counted every citation of both IDs and sorted them by the half each one means, which is the evidence that decides between A and B, read the matrix's ID pattern to rule out C, and chose A.
It also found the unlogged \texttt{num\_\allowbreak{}threads} while writing NFR-14's route, and its first recommendation was to leave that and the \texttt{tracking.\allowbreak{}py} comments to a later change that retrains anyway.
An independent AI review of the pull request rejected that, because it left NFR-14 reported as verified by a test while its parameter clause was unmet, and the fix was made in this pull request, test first, which is the modification.},
  aievidence = {Claude Code session on 2026-10-05 implementing issue \#60: \texttt{grep -\allowbreak{}rn \textquotedbl{}NFR-\allowbreak{}0[69]\textquotedbl{}} over the documents, the report, the tests, the tools and the package, each hit classified by the half it means, and \texttt{git grep} counts on \texttt{main} for the figures above; the requirement matrix rebuilt after the split, with NFR-14 verified by five tests and NFR-15 verified by hand; \texttt{test\_\allowbreak{}a\_\allowbreak{}run\_\allowbreak{}records\_\allowbreak{}every\_\allowbreak{}parameter\_\allowbreak{}its\_\allowbreak{}stage\_\allowbreak{}declares} run against the unchanged \texttt{train.\allowbreak{}py} first, failing for all four variants on \texttt{num\_\allowbreak{}threads} alone.},
  evidence = {\href{https://github.com/mlops-2627q1-mds-upc/MLOPS_Recommenditos/issues/60}{issue \#60} and its comment on NFR-06; \href{https://github.com/mlops-2627q1-mds-upc/MLOPS_Recommenditos/pull/70}{PR \#70}; \href{https://github.com/mlops-2627q1-mds-upc/MLOPS_Recommenditos/blob/main/docs/docs/requirements.md}{requirements} and \href{https://github.com/mlops-2627q1-mds-upc/MLOPS_Recommenditos/blob/main/docs/docs/specification.md}{specification} NFR-06, NFR-09, NFR-14 and NFR-15; \texttt{tools/\allowbreak{}expected\_\allowbreak{}coverage.\allowbreak{}yaml}; \texttt{DUE\_\allowbreak{}AT\_\allowbreak{}M3} in \texttt{tests/\allowbreak{}test\_\allowbreak{}requirement\_\allowbreak{}matrix.\allowbreak{}py}; \hyperref[EDN-27]{EDN-27} for the one ID space, \hyperref[EDN-44]{EDN-44}, \hyperref[EDN-45]{EDN-45}, \hyperref[EDN-46]{EDN-46}, \hyperref[EDN-56]{EDN-56}, \hyperref[EDN-71]{EDN-71}.},
}
